\documentclass[a4paper,14pt]{extarticle}
\usepackage{styles}

\begin{document}

\litSubsubsection{2.1.2 CUDA Model Suitability for GA-Specific Computations}

\textbf{CUDA Model Suitability for GA-Specific Computations.} Compute Unified Device Architecture - ``CUDA is a parallel computing platform and programming model developed by NVIDIA that enables dramatic increases in computing performance by harnessing the power of the GPU''~\cite{nvidia2023cuda}. CUDA organizes computation into kernels executed by lightweight threads arranged in blocks and grids, with access to a hierarchical memory system (global, shared, constant, and texture memories). This architecture is well suited to data-parallel problems where the same operation is applied independently across large sets of data - a pattern that matches the population-based structure of GAs, where fitness evaluation and variation operators are applied repeatedly across individuals. The CUDA platform exposes GPU architectural features directly to programmers, enabling fine-grained control over thread execution to fully exploit GPU computing power~\cite{cheng2020parallel}.

\textbf{Python-CUDA Integration for GA Development.} The primary CUDA API uses C, but Python interfaces exist: CuPy provides NumPy-compatible GPU arrays backed by cuBLAS and cuFFT~\cite{cupy_docs}; Numba provides a JIT compiler that generates CUDA kernels directly from Python via the \texttt{@cuda.jit} decorator~\cite{numba_cuda}. These tools allow GA practitioners to implement operations in Python and delegate numeric workload to CUDA kernels on the GPU. ``Numba exposes the CUDA programming model, just like in CUDA C/C++, but using pure python syntax, so that programmers can create custom, tuned parallel kernels without leaving the comforts and advantages of Python behind.''\cite{harris2017numba}

\textbf{Benefits and Drawbacks of Python for CUDA GA Development.} Python is chosen for CUDA development in this paper to balance performance with rapid prototyping and accessibility. CuPy compiles custom kernels at runtime using NVRTC (NVIDIA Runtime Compilation), producing PTX that the driver compiles to device code, with results cached for reuse~\cite{cupy_docs}. Numba's JIT achieves performance close to native CUDA: ``CUDA Python Mandelbrot code runs nearly 1700 times faster than the pure Python version''~\cite{harris2017numba}. Python's concise syntax and ecosystem make it effective for disseminating GPU-accelerated GA techniques beyond systems programming specialists.

First, Python's interpreter adds host-side overhead to every GPU call; for small inputs the launch overhead can exceed the computation itself~\cite{carpentries_gpu}. Second, the Python-CUDA stack depends on multiple native libraries whose versions must align: CuPy requires a matched CUDA toolkit wheel~\cite{cupy_install}, and Numba requires both the toolkit and a compatible driver~\cite{numba_cuda}. Third, Python's GIL limits multi-core CPU orchestration of multiple GPUs, as it ``prevents multiple threads from executing Python code at the same time''~\cite{pep703}. These drawbacks are mitigated by focusing on large, compute-bound GA workloads with most computation on the device.

\end{document}
